%% ============================================================
%%  LegalShadow — Cover letter for the REVISED submission
%%  Author's Reply to the Academic Editor · MDPI SuSy
%%  Engineering Proceedings (ISSN 2673-4591) · ICI2ST 2026 volume
%%  Standalone: compile with pdflatex. No external class required.
%%  AUTOR: completar el Manuscript ID de SuSy antes de enviar.
%% ============================================================
\documentclass[11pt,a4paper]{article}

\usepackage[T1]{fontenc}
\usepackage[utf8]{inputenc}
\usepackage{lmodern}
\usepackage[english]{babel}
\usepackage{amsmath}
\usepackage[left=2.4cm,right=2.4cm,top=2.4cm,bottom=2.4cm]{geometry}
\usepackage{microtype}
\usepackage{parskip}
\usepackage{enumitem}
\usepackage{xcolor}
\usepackage[hidelinks]{hyperref}

\definecolor{hi}{HTML}{B45309}
\setlist[itemize]{leftmargin=1.4em, itemsep=3pt, topsep=4pt}
\setlist[enumerate]{leftmargin=1.6em, itemsep=3pt, topsep=4pt}
\newcommand{\Psiglob}{$\Psi_{\mathrm{global}}$}
\pagestyle{empty}

\begin{document}

%% ── Sender ──────────────────────────────────────────────────
{\raggedright
\textbf{Patricio M. Paccha-Angamarca}\\
Software Engineering Creation, Application and Management Group (GI-CreAGSoft)\\
National Polytechnic School, Ladrón de Guevara E11-253, Quito 170525, Ecuador\\
\href{mailto:patricio.paccha@epn.edu.ec}{patricio.paccha@epn.edu.ec} \quad ORCID: 0000-0001-6285-7390\\[6pt]
Quito, Ecuador --- \today\par}

\vspace{10pt}

%% ── Addressee ───────────────────────────────────────────────
{\raggedright
\textbf{To the Academic Editor}\\
\emph{Engineering Proceedings} (MDPI, ISSN 2673-4591)\\
ICI2ST 2026 --- International Conference on Information Systems and Software Technologies\\
% AUTOR: insertar el identificador que asigna SuSy.
Manuscript ID: engproc-XXXXXXX\par}

\vspace{12pt}

\textbf{Subject:} Submission of the revised manuscript and author's reply to the Academic Editor

\vspace{8pt}

Dear Academic Editor,

Thank you for a report that was demanding in the most useful way. All five major points identified real weaknesses, and addressing them changed the paper for the better: the delimitation requested in §3.3 forced us to state the \emph{direction} of validation precisely, and that turned out to be the clearest expression of what the contribution actually is. We enclose the revised manuscript together with a detailed point-by-point response.

%% ── Title change ────────────────────────────────────────────
\vspace{4pt}
{\color{hi}\textbf{Change of title (one word).}}
We wish to draw your attention to a single change in the title, which we make explicit here so that it is not mistaken for an inconsistency between versions:

\vspace{4pt}
\begin{quote}
\emph{Previous:} LegalShadow: An Ontology-Guided Digital Shadow Architecture for Semantic Validation of \textbf{Judicial Digital} Ecosystems\\[4pt]
\emph{Revised:} LegalShadow: An Ontology-Guided Digital Shadow Architecture for Semantic Validation of \textbf{LegalTech} Ecosystems
\end{quote}

\vspace{2pt}
The revised wording places the paper within the LegalTech research line to which it belongs and aligns the title with the governance ontology that the instrument validates against, in which the LegalTech domain is one of the ten dimensions. No other element of the title, and no claim in the paper, changes as a consequence.

%% ── Summary of revisions ────────────────────────────────────
\vspace{6pt}
\textbf{Summary of the revisions.}
Every change is marked in the \LaTeX{} source with \verb|%>> EDITOR| comments, so each requested item can be located directly in the manuscript.

\begin{itemize}
  \item \textbf{§3.1 External validity.} A dedicated block now separates three claims the previous version allowed a reader to conflate. The \emph{numerical} result generalises to nothing---46/100 and \Psiglob{}~$=34.9\%$ describe one institution over a one-month window of 17 dated runs. The \emph{diagnostic pattern} is stated as a hypothesis that only a multi-institutional study can test. Only the \emph{instrument} generalises, and we give a four-step transfer protocol that names the one step carrying real cost: roughly a fifth of the anchoring rules key on national artefacts and must be re-expressed against equivalent infrastructure elsewhere.
  \item \textbf{§3.2 Scoring and the mapping $\Gamma$.} A new subsection states what $\Gamma$ actually is---a versioned, deterministic table of \emph{signal $\rightarrow$ class} rules, each firing on a syntactic condition re-evaluable on the stored snapshot---with concrete instances. We also clarified something the previous version left implicit and which we believe caused the opacity you identified: the ontology \emph{drives} the acquisition, so $\Gamma$ is that same rule table read in the opposite direction. Separately, we now distinguish the two quantities by epistemic status: $\mathrm{onto}(d)$ is expert-elicited, whereas $\Psi(d)$ is fully deterministic and reproducible from the published artefacts. This makes partial replication possible for a reviewer who rejects our rubric.
  \item \textbf{§3.3 Novelty.} A new Section~2, \emph{Related Work and Delimitation}, with nine new third-party references and a comparison table, positions the work against the three lines you name. The claim is no longer asserted but bounded to the scope of that comparison, and the phrase ``to the best of our knowledge'' has been removed.
  \item \textbf{§3.4 Presentation.} A JSON-LD listing shows one real shadow component with its evidence block, anchoring array and recorded gap. A worked example carries the \emph{security posture} dimension end to end, from HTTP evidence to its position on the remediation agenda. The four-versus-ten contrast is now a titled paragraph of its own.
  \item \textbf{§3.5 Reflexive finding.} Developed along the two axes you propose, which pull in opposite directions: a methodological \emph{measurement floor}---coverage is a joint property of the ecosystem and of its publication practices, so a low $\Psi$ cannot distinguish weak governance from undisclosed governance---and a concrete publication agenda of three low-cost acts, none of which requires reforming how the institution is governed.
  \item \textbf{§4 Minor issues.} Over-claims removed, terminology fixed on first mention, every figure and table now described in the running text, and the relationship with our previous MALTG work stated as an explicit inventory of what is reused unchanged and what is new.
\end{itemize}

%% ── Points flagged proactively ──────────────────────────────
\vspace{4pt}
\textbf{Two points we flag proactively.}
First, our response to §3.1 deliberately \emph{weakens} several claims about generalisation rather than strengthening them. We judged that a narrow claim supported by the evidence serves the paper better than a broad one a second jurisdiction might contradict, but we will revisit this balance if you read it differently. Second, the revised Discussion was consolidated from eight short paragraphs into five, because three of them restated the reflexive finding and transferability. No content was removed in doing so; the section is denser, not shorter.

\vspace{4pt}
All artefacts required to reproduce every figure---source code, the OWL governance ontology, the versioned scoring rubric, the anchoring rule table, the dated JSON-LD shadows and the SHA-256 evidence manifests---are openly deposited at Zenodo (DOI:~10.5281/zenodo.21942315). We would welcome reviewers exercising the artefact directly: the environment builds with a single \texttt{docker compose} invocation.

\vspace{4pt}
We hope the revision meets the standard for a second round, and we thank you again for a report that improved the work.

\vspace{14pt}
{\raggedright
Yours sincerely,\\[20pt]
\textbf{Patricio M. Paccha-Angamarca}\\
\emph{on behalf of the co-authors} Erwin J. Sacoto-Cabrera and Víctor V. Velepucha-Bonett\\
Corresponding author --- \href{mailto:patricio.paccha@epn.edu.ec}{patricio.paccha@epn.edu.ec}\par}

\end{document}
